%!TEX TS-program = lualatex
\documentclass[
    language=ukrainian,
    doctype=article,
]{../../omnilatex}

\title{Штучний інтелект у вищій освіті}
\subtitle{Виклики, можливості та перспективи розвитку}
\author{Д-р Олександр Петренко}
\date{\today}
\keywords{штучний інтелект, вища освіта, машинне навчання, цифрова освіта}

\begin{document}
\maketitle

\begin{abstract}
Штучний інтелект глибоко трансформує методи навчання та викладання у вищих навчальних закладах. Ця стаття досліджує основні сфери застосування штучного інтелекту у вищій освіті, включаючи персоналізоване навчання, автоматизоване оцінювання та аналітику прогнозування академічної успішності студентів. Ми також обговорюємо етичні та інституційні виклики, пов'язані з цією трансформацією.
\end{abstract}

\section{Вступ}
Вища освіта стикається з небаченим обсягом даних. Щороку генеруються трильйони точок даних з медичних зображень, секвенування геному, електронних медичних карт та клінічних випробувань. Штучний інтелект та методи машинного навчання пропонують нові підходи до використання цих даних для відкриття механізмів захворювань, розробки діагностичних інструментів та створення нових терапій.

\section{Персоналізоване навчання}
Системи персоналізованого навчання використовують штучний інтелект для адаптації освітнього досвіду до індивідуальних потреб кожного студента. Ці системи аналізують стилі навчання, сильні та слабкі сторони студентів, а потім пропонують найбільш підходящий навчальний матеріал та завдання.

Дослідження показують, що системи персоналізованого навчання можуть значно покращити академічні результати порівняно з традиційними методами. Студенти, які користуються такими системами, як правило, отримують вищі оцінки, мають нижчий рівень відсіву та вищу задоволеність навчанням.

\section{Автоматизоване оцінювання}
Штучний інтелект дозволяє швидко та послідовно оцінювати студентів. Системи автоматичного оцінювання тексту можуть надійно оцінювати есе, тоді як системи аналізу коду можуть перевіряти та оцінювати програмні завдання.

Основна перевага автоматизованого оцінювання - це можливість надати студентам негайний зворотний зв'язок, що прискорює процес навчання. Крім це, автоматизація знижує навантаження на викладачів, дозволяючи їм зосередитися на індивідуальному підході до студентів.

\section{Аналітика прогнозування}
Прогнозна аналітика використовує минулі дані для прогнозування майбутніх результатів. У вищій освіті цей підхід може бути використаний для визначення студентів, які мають ризик припинити навчання, виявлення тих, хто потребує додаткової підтримки, та покращення стратегій набору.

Багато навчальних закладів використовують прогнозну аналітику для визначення ризикованих студентів на ранньому етапі та втручання до того, як буде пізно. Цей підхід показує перспективні результати у покращенні рівня завершення навчання.

\section{Етичні виклики}
Використання штучного інтелекту у вищій освіті створює кілька етичних викликів. Перше - це конфіденційність даних. Дані про навчання студентів є особистою та чутливою інформацією, і використання цих даних для навчання моделей штучного інтелекту вимагає згоди студентів.

Друге - це справедливість алгоритмів. Якщо дані для навчання не представляють усі групи студентів, моделі можуть демонструвати систематичну упередженість. Це може призвести до дискримінації студентів з недостатньо представлених груп.

\section{Висновки}
Штучний інтелект має величезний потенціал для трансформації вищої освіти. Однак успіх цієї технології залежить від обережного застосування з урахуванням етичних принципів, конфіденційності та справедливості. Співпраця між вченими в галузі комп'ютерних наук, викладачами та регуляторними органами є необхідною для повного розкриття потенціалу штучного інтелекту в освіті.

\end{document}
